% Proof note for the positive-D regional completion.
% Reviewed mathematical rewrite; regional execution completed.
% All 55 original lower coefficient identities, 10 certificate aggregates,
% and 10 positivity consumers passed. CompleteRow2, CheckpointSeventeen,
% and the LeanVerification root compiled successfully on 25 November 2026.
% The final combined build, transitive axiom audit, and release review
% are recorded separately in verification_status.json.

\subsection{Positive $D$: a signed real secant proof}
\label{lean17:positive-d-proof}

The full regional statement is
\[
 \frac14\le F<1,\quad D>0
 \quad\Longrightarrow\quad
 \operatorname{CriticalPeriodBound}(1,D,F).
\]
The formal bound counts distinct isolated critical energies of the least
period of the original Loud vector field, on its full period annulus and
in the original time variable. The interior argument uses the actual two
real turning points. Its domain is
\[
 D>0,\qquad \frac14<F<1,\qquad F\ne\frac12,\qquad
 h\in\operatorname{AdmissibleEnergy}(D,F).
\]
The quarter edge has the established bound zero. The half line is treated
by continuity of physical width derivatives after the interior comparison.

\paragraph{Normalization and the signed branch.}
Put $a=F^{-1}-1$ and $d=-(D+1)/D$. Then $0<a<3$, $a\ne1$, and
$d<-1$. The mechanical potential is an affine energy rescaling of
\[
 V(u)=\frac d2u^2+\frac{1-d}{1-a}u^{1-a}
       +\frac1{2a}u^{-2a}.
\]
Write the actual equal-energy turns as
$u_+=\nu e^q$, $u_-=\nu e^{-q}$, with $\nu,q>0$, and set
\[
 f=\frac{\tanh(aq)}{\tanh q},\qquad
 \rho=\left(\frac{\sinh(aq)}{a\sinh q}\right)^2,\qquad
 r=\nu^{1+a}\frac{\sinh q}{\sinh(aq)},\qquad
 \omega=ad r^2-1.
\]
Thus $\omega<-1$. The equal-energy equation gives
\[
 f(1+\omega)+\frac{a(1-d)}{1-a}(1-f)r-1=0.
\]
The force signs at the two actual turns prove
$1+f(1+\omega)<0$. The relevant identity is
\[
 a\nu^{2a}\{u_+V'(u_+)-u_-V'(u_-)\}
 =(1+a)\{ad\nu^{2(1+a)}\sinh(2q)+\sinh(2aq)\},
\]
whose right-hand bracket has the sign of $1+f(1+\omega)$.
The files \texttt{SignedQuadratic}, \texttt{PhysicalSheet}, and
\texttt{SignedTransport} establish this branch and its derivatives.

\paragraph{Regularity and signed denominators.}
Write
\[
 g=1-a^2\rho,\qquad A=f^2-a^2\rho,\qquad
 C=a\rho(1-af),\qquad M_0=g(1+f),\qquad N=M_0+C\omega.
\]
Here $M_0$ is distinct from the auxiliary polynomial
$\texttt{sourceM}=(1+a)(f-a\rho)$. The secant geometry supplies the
seven exclusions recorded in \texttt{SignedRegular}:
\[
 a\ne0,\quad 1+f\ne0,\quad g\ne0,\quad f-a\ne0,\quad
 1-f\ne0,\quad 1-\rho\ne0,\quad A\ne0.
\]
The actual turning-point transport gives $Ng<0$; the secant geometry
gives $Ag>0$ and $Cg>0$. In particular $C,N\ne0$, $AN<0$, and
\[
 n=-\frac NC>0,\qquad
 \Delta_P=4(1+\omega)AN^3>0,\qquad
 \Delta_R=4a(1+f)(1+\omega)AN^5>0.
\]
Indeed $(1+\omega)AN>0$, and the remaining powers of $N$ are even.
The coefficient $C$ has opposite signs on the two secant cones; these
denominator statements hold on both. The named sign lemma is
\texttt{signed\_numerator\_domain}, applied to the actual turns by
\texttt{actual\_signed\_numerator\_domain}.

\paragraph{Energy and width jets.}
Throughout this paragraph $D,F$ are fixed, and primes denote derivatives
with respect to the original physical energy $h$. The normalized energy
is affine:
\[
 \widetilde E(h)=\frac{h-h_0}{c},\qquad
 h_0=\operatorname{zeroEndpointHeight}(D,F),\qquad
 c=\frac{D+1}{F}>0.
\]
Here $c$ is exactly the product of the physical energy scale and $d$.
Hence $\widetilde E'=1/c$ is constant and $\widetilde E''=0$.
Let $W=u_+-u_-=F W_{\rm phys}$ and $\lambda=W'/W>0$.
The exact physical energy-width identity is
\[
 \widetilde E=W^2U,\qquad
 U=\frac{(1+f)^2\{(1+\omega)^{-1}-a^2\rho\}}{8A},
 \qquad \widetilde E'>0.
\]
The rational transport produces $k,\Phi$ with
\[
 U'=\lambda(k-2U),\qquad k'=-\lambda\Phi,\qquad
 \widetilde E'=WW'k.
\]
Consequently $k>0$. Differentiating the last identity and using
$\widetilde E''=0$ and $k'=-(W'/W)\Phi$ gives
\[
 0=\{(W')^2+WW''\}k+WW'k'
  =\{(W')^2+WW''\}k-(W')^2\Phi.
\]
On $\Phi>0$, the checked diagonal jets and the signed regularity
conditions give
\[
 \Phi'=\lambda(R_X-\Phi),\qquad
 R_X=\operatorname{modelF}\,\operatorname{modelR}.
\]
The finite estimate needed for the physical argument is therefore
$\Phi>0\Longrightarrow R_X>0$.
The affine normalization and transport identities are in
\texttt{SignedNormalization}, \texttt{PhysicalEnergyRatio},
\texttt{SignedEnergy}, \texttt{PhysicalModel}, and
\texttt{ActualCurvatureJet}.

\paragraph{The two strict secant cones.}
For $0<a<1$ one has
\[
 a<f<1,\qquad 0<\rho<1,\qquad f\rho<a.
\]
For $1<a<3$ one has
\[
 1<f<a,\qquad \frac af<\rho<\mathcal R(a,f),
\]
where
\[
 \begin{split}
 B_{\rm env}(a,f)={}&-a^3f+2a^2f^2+3a^2f+2a^2
                    -4af^2-3af-4a+5f,\\
 \mathcal R(a,f)={}&
 \frac{f^2(a-1)^3(af^2+a-2f)}{a^2(f-1)^2B_{\rm env}(a,f)}.
 \end{split}
\]
The denominator $B_{\rm env}$ is positive on the lower strip.
The lower correlation follows from
\[
 H=(a-1)^3f^3(f+a^2\rho)
   -a^2\rho(f-1)^2\{a^2\rho(f^2+1)-2f^2\}>0
\]
and the denominator-free tangent identity
\[
 \begin{split}
 f^2\{f^2(a-1)^3(af^2+a-2f)-a^2(f-1)^2B_{\rm env}\rho\}
  ={}&(f-1)^2H\\
    &+(f^2+1)\{a^2(f-1)^2\rho-(a-1)^2f^2\}^2.
 \end{split}
\]
Both correlations use strict monotonicity of
$\sinh^3q/(q^3\cosh q)$. The lower comparison uses
$(a-1)q<2q$, which retains the restriction $a<3$.
These estimates are in \texttt{HyperbolicBounds}, \texttt{LowerCone},
and \texttt{SecantGeometry}.

\paragraph{Positive forms of the lower coordinate denominators.}
Put
\[
 f_c=1+(a-1)t,\qquad
 C_c=2+2(a-1)t+a(a-1)t^2,\qquad
 D_c=a^2t^2B_c(a,t),\qquad \mathcal D=f_cD_c.
\]
The low-degree coordinate polynomials have the exact positive forms
\[
 \begin{split}
 B_c(a,t)={}&(5-a)(1-t)^2+(5-a)(a+1)t(1-t)\\
            &+a(a+1)t^2,
 \end{split}
\]
and
\[
 \begin{split}
 G_c(a,t)={}&a(a-1)^4t^3+
       \{a(a-1)((a-2)^2+5)+2\}t^2\\
            &+(8a-4)t+2.
 \end{split}
\]
They give $B_c>0$ for $0<a<5$, $0<t<1$, and $G_c\ge2$ for
$a\ge1$, $t\ge0$. In the required lower cone $1<a<3$ and
$0<t<1$, all of $f_c,D_c,\mathcal D$ are strictly positive.
The wider algebraic ranges of these two lemmas leave the hyperbolic
domain $a<3$ unchanged.

The coordinate identities are
\[
 B_{\rm env}(a,f_c)=(a-1)^2B_c(a,t),\qquad
 f_c^3C_c-aD_c=(1-t)^2G_c(a,t).
\]
They expose the gap in the correlated interval:
\[
 \mathcal R(a,f_c)-\frac a{f_c}
 =\frac{(1-t)^2G_c(a,t)}{\mathcal D}>0.
\]
With $a=1+2x$, the lower chart consequently takes the form
\[
 \rho=\frac{aD_c+(1-t)^2G_c(a,t)\eta}{f_cD_c},
 \qquad (x,t,\eta)\in(0,1)^3.
\]
The exact identities and their sign proofs are recorded in
\texttt{LowerCoordinates} and \texttt{LowerCoordinateStructure}.

\paragraph{Rational normalization and homogeneous pole coordinates.}
The source normalization first cancels independent denominator factors;
the source coefficient identities are separate integer polynomial
equalities. Throughout the regular sheet the first fraction below holds.
For the implication under consideration, assume $\Phi>0$.
This derives the required $\texttt{modelPsi}\ne0$ and hence the second:
\[
 \Phi=\frac{P_\omega(\omega)}{\Delta_P},\qquad
 R_X=\frac{R_\omega(\omega)}{\Delta_R},\qquad
 \deg P_\omega\le4,\quad \deg R_\omega\le6.
\]
For a coefficient expansion
$Q_\omega(s)=\sum_{i=0}^m Q_i s^i$, with $m=4$ or $m=6$, define
\[
 \widehat Q_j=\sum_{i=j}^m
       \binom ij(-1)^iM_0^{\,i-j}C^{m-i}Q_i.
\]
The homogeneous change of variables is the polynomial identity
\[
 C^mQ_\omega(\omega)=\sum_{j=0}^m\widehat Q_j(-N)^j.
\]
Set $p_j=C^j\widehat P_j$, $r_j=C^j\widehat R_j$, and define
\[
 P_n(s)=\sum_{j=0}^4p_js^j,\qquad
 R_n(s)=\sum_{j=0}^6r_js^j.
\]
Since $C\ne0$ and $n=-N/C$,
\[
 P_n(n)=C^4P_\omega(\omega),\qquad
 R_n(n)=C^6R_\omega(\omega).
\]
Both $C^4$ and $C^6$ are strictly positive. Thus this change preserves
the signs needed in the comparison on both secant cones. The formal
identities are \texttt{inverse\_homogeneous\_4},
\texttt{inverse\_homogeneous\_6}, \texttt{phiPoleValue\_actual}, and
\texttt{rxPoleValue\_actual}.

\paragraph{One quartic--sextic comparison and ten scalar leaves.}
For $n>0$, the coefficient conditions used by the abstract comparison are
\[
 p_4<0,\quad r_0\ge0,\quad r_1>0,\quad r_2,r_3,r_4,r_5\ge0,
 \qquad
 s_j=-p_4r_{j+2}+r_6p_j\ge0\quad(0\le j\le3).
\]
When $r_6\ge0$, these signs give $R_n(n)>0$ directly.
When $r_6<0$, use
\[
 -p_4R_n(n)=-p_4r_0-p_4r_1n+
             \sum_{j=0}^3s_jn^{j+2}-r_6n^2P_n(n).
\]
If $P_n(n)>0$, the right side is positive; in particular
$-p_4r_1n>0$. This is
\texttt{sextic\_pos\_of\_quartic\_pos}. Combining it with the positive
denominators and the positive powers of $C$ gives the required
$\Phi>0\Longrightarrow R_X>0$.

The ten finite leaves have a fixed role: leaf 00 controls $-p_4$;
leaves 01--05 control $r_1,\ldots,r_5$; and leaves 06--09 control
$s_0,\ldots,s_3$. The remaining constant coefficient factors as
\[
 r_0=
 3a^2\rho^3(1+a)^6(1+f)^5g^3(f-a\rho)^4K_0^2(1-f),
\]
where $K_0=\texttt{poleCore00}(a,f,\rho)$. Its nonnegativity follows
from $g(1-f)>0$ and the even powers, including the square $K_0^2$.
The structure \texttt{LeafSigns} uses sign $+1$ on $a<1$ and $-1$
on $a>1$; only leaves 01 and 05 carry this sign. The conversion lemmas
\texttt{upper\_pole\_coefficient\_signs} and
\texttt{lower\_pole\_coefficient\_signs} supply exactly the
\texttt{PoleCoefficientSigns} needed by the abstract comparison.

\paragraph{Finite positivity on the two cubes.}
The upper chart is $f=a+(1-a)t$. Its first eight leaves use
$\rho=\eta$ on $(a,t,\eta)\in(0,1)^3$.
Leaves 08 and 09 use $\rho=a\eta/f$; the physical correlation
$f\rho<a$ is exactly what places $\eta$ in $(0,1)$.
The upper certificates list $8584$ positive integer Bernstein
coefficients.

The lower chart keeps the full envelope correlation through the
positive denominator $\mathcal D$. Write $W_c=(1-t)^2G_c(a,t)$.
For $L(\rho)=\sum_{i=0}^m c_i\rho^i$ the exact homogeneous formula is
\[
 \begin{split}
 &(f_cD_c)^m
 L\left(\frac{aD_c+W_c\eta}{f_cD_c}\right)\\
 &\quad=
 \sum_{j=0}^mD_c^{m-j}W_c^j
       \left(\sum_{i=j}^m\binom ij c_i a^{i-j}f_c^{m-i}\right)\eta^j.
 \end{split}
\]
The formal homogeneous lemmas cover degrees $2,3,4,6$.
The inner sums remain named bivariate coordinate cores.
Boundary monomial factors are extracted explicitly and are positive
on the open cube. Each remaining atom is identified exactly with a
finite sum of unnormalized Bernstein terms
\[
 c_{ijk}x^i(1-x)^{d_x-i}
         t^j(1-t)^{d_t-j}
         \eta^k(1-\eta)^{d_\eta-k}.
\]
Every listed coefficient is a positive integer, all indices satisfy
their degree bounds, and the list is nonempty. Omitted coefficients
are zero. The lower certificates list $59167$ positive coefficients.
These statements concern the Bernstein representation; positivity
of the coefficients in an expanded monomial basis is not required.

All coordinate denominators are proved positive before their powers
are cleared. The finite obligations are exact polynomial identities
and finite positivity proofs. Separate generated modules retain the
checked pieces and assemble the same source identities; none of
their validity or source-correspondence obligations is added as an
assumption of the regional theorem.

\paragraph{The checked source-identity interface.}
For each large coefficient identity, Lean checks four propositions:
validity of all certificate operations, equality of the two selected
output caches, and two source-erasure equalities identifying the two
output expressions with \texttt{lowerPowers} of the unchanged source expressions.
The theorem \texttt{CertificateKernel.sourceIdentity\_of\_checked} combines
these four proved statements with the evaluation semantics to give equality
of the two \texttt{SourceValue} evaluations. The original real-formula
wrapper then gives its unchanged real identity. Independent \texttt{Guards}
and \texttt{Erase} modules use shared transparent data; a short
\texttt{Kernel} module composes their conclusions.

\paragraph{Return to the actual period and the seams.}
The width identities and finite comparison imply
\[
 (W')^2+WW''>0\Longrightarrow W'W'''-3(W'')^2>0.
\]
The positive scaling $W=F W_{\rm phys}$ transports this implication
to the physical width. Set
\[
 Q=-W_{\rm phys}/4+hW_{\rm phys}'+h^2W_{\rm phys}'',\qquad
 \chi=(W_{\rm phys}')^2+W_{\rm phys}W_{\rm phys}''.
\]
The physical width criterion proves $Q>0\Longrightarrow\chi>0$.
It then uses the exact identity
\[
 4W_{\rm phys}'\,hQ'
 =3h\{W_{\rm phys}'+2hW_{\rm phys}''\}^2
  +4h^3\{W_{\rm phys}'W_{\rm phys}'''-3(W_{\rm phys}'')^2\}.
\]
Thus $hQ'+Q/2>0$ wherever $Q>0$.
For every fixed admissible cap $H$, the physical Abel input has a
continuous extension on $[0,H]$ given by
\[
 z(0)=0,\qquad z(h)=\sqrt h\,Q(h)\quad(0<h\le H).
\]
With the original time normalization,
\[
 (Az)(h)=2\int_0^{\pi/2}z(h\sin^2\theta)\,d\theta,\qquad
 \sqrt2\,(Az)(h)=h^2T'(h).
\]
The checked Abel-period transfer gives the bound one for the original
least period $T$. It also covers the alternative in which the Abel
input vanishes throughout the annulus: then $T'$ vanishes throughout,
and there are no isolated critical energies.

At $F=1/2$, fix an admissible physical energy with $\chi>0$.
Joint analyticity and openness of admissibility preserve these
conditions for nearby $F>1/2$. Continuity transfers the nonnegative
physical Schwarz expression to the seam. The strict Euler conclusion
still follows from $Q>0$ and the displayed width identity.
The $F=1/4$ edge uses the already established bound zero.
These two endpoint arguments close the full parameter region.

The final export is
\texttt{PositiveDCompletion.full\_regional\_row2} in
\texttt{CompleteRow2.lean}; its only mathematical premise is
\texttt{parameterRegion 2 D F}. The global assembly combines this
regional conclusion with the other table bounds and the analytic
sharpness construction. The latter supplies ordered simple critical
energy paths, with negative second period derivative at the first
and positive second derivative at the second. The final build and transitive
axiom audit of these exact exports are recorded in
\texttt{verification\_status.json}.
